
\documentclass[12pt, a4paper, oneside]{ctexart}
\usepackage{amsmath, amsthm, amssymb, graphicx}
\usepackage[bookmarks=true, colorlinks, citecolor=blue, linkcolor=black]{hyperref}


\title{高等代数与解析几何1期末考试题}
\author{}
\date{2012.7.20}

\begin{document}

\maketitle
\;
\\一，计算下列行列式：
\\
(1)\\
$$\begin{vmatrix}0 & 2 & 3 & \cdots & n\\
			 1 & 0 & 3 & \cdots & n\\
			 1 & 2 & 0 & \cdots & n\\
			 \vdots & \vdots & \vdots & \; & \vdots \\
			 1 & 2 & 3 & \cdots & 0\\\end{vmatrix}$$
(2)\\
$$\begin{vmatrix}
x_1 & y & y & \cdots & y\\
-y & x_2 & y & \cdots & y\\
-y & -y & x_3 & \cdots & y\\
\vdots & \vdots & \vdots & \; & \vdots\\
-y & -y & -y & \cdots & x_n
\end{vmatrix}$$
(3)\\
$$\begin{vmatrix}
1+x_1y_1 & 1+x_1y_2 & 1+x_1y_3 & \cdots & 1+x_1y_n\\
2+x_2y_1 & 2+x_2y_2 & 2+x_2y_3 & \cdots & 2+x_2y_n\\
3+x_3y_1 & 3+x_3y_2 & 3+x_3y_3 & \cdots & 3+x_3y_n\\
\vdots & \vdots & \vdots & \; & \vdots\\
n+x_ny_1 & n+x_ny_2 & n+x_ny_3 & \cdots & n+x_ny_n
\end{vmatrix}$$

\;
\\二，求当a,b取何值时该线性方程组有无穷多组解，并写出此时的结构通解.
$$\begin{cases} 
ax_1+x_2+x_3+x_4=1\\
x_1+ax_2+x_3+x_4=b\\
x_1+x_2+ax_3+x_4=b^2\\
x_1+x_2+x_3+ax_4=b^3
\end{cases}$$

\;
\\三，求该向量组的一个极大无关组，并将其余向量用这个极大无关组中的向量表示.\\

\;

$\alpha_1$=$\left (\begin{array}{rrrr}
1\\
-1 \\
2 \\
4
\end{array}\right)$
$\alpha_2$=$\left (\begin{array}{rrrr}
0\\
3 \\
1 \\
2
\end{array}\right)$
$\alpha_3$=$\left (\begin{array}{rrrr}
3\\
0 \\
7 \\
14
\end{array}\right)$
$\alpha_4$=$\left (\begin{array}{rrrr}
1\\
-1 \\
2 \\
0
\end{array}\right)$
$\alpha_5$=$\left (\begin{array}{rrrr}
3\\
0 \\
7 \\
6
\end{array}\right)$

\;

\;
\\四，求下列矩阵的逆矩阵：
\;
\\(1)\\
$$\left (\begin{array}{rrrrr}
0 & 1 & 1 & \cdots & 1\\
1 & 0 & 1 & \cdots & 1\\
1 & 1 & 0 & \cdots & 1\\
\vdots & \vdots & \vdots & \; & \vdots\\
1 & 1 & 1 & \cdots & 0
\end{array}\right)$$

\;
\\(2)\\
$$\left (\begin{array}{rrrr}
6 & 5 & 11 & 0\\
5 & 4 & 0 & 11\\
0 & 0 & 4 & -5 \\
0 & 0 & -5 & 0
\end{array}\right)$$
\;\\
(3)\\$$\mathrm{I}_n-\alpha\beta^{\mathrm{T}}$$
其中$\beta^{\mathrm{T}}\alpha\neq1$，$\alpha,\beta\in\mathrm{R}^n$.

\;\\
五，设$\mathrm{A},\mathrm{B}$是n级矩阵，$\mathrm{A}\mathrm{B}=\mathrm{A}+2\mathrm{B}$\\
(1)证明：$\mathrm{A}\mathrm{B}=\mathrm{B}\mathrm{A}$\\
(2)求$(2\mathrm{I}_n-A)^{-1}$

\;\\
六，设$\mathrm{A}$是$m\times n$矩阵，$\mathrm{B}$是$n\times m$矩阵，证明：\\
(1)$\mathrm{det}(\mathrm{I}_n-\mathrm{B}\mathrm{A})=\mathrm{det}(\mathrm{I}_m-\mathrm{A}\mathrm{B})$\\
(2)$\mathrm{rank}(\mathrm{I}_n-\mathrm{B}\mathrm{A})=\mathrm{rank}(\mathrm{I}_m-\mathrm{A}\mathrm{B})$

\;\\
七,证明$\{1,e^x,\dots e^{x^n}\}$线性无关.

\;\\
八，证明
$$\left (\begin{array}{rrrrrr}
1 & -2 & 0 & \cdots & 0 & 0\\
2 & 1 & -2 & \cdots & 0 & 0\\
0 & 2 & 1 & \cdots & 0 & 0\\
\vdots & \vdots & \vdots & \; & \vdots &\vdots\\
0 & 0 & 0 & \cdots & 1 & -2\\
0 & 0 & 0 & \cdots & 2 & 1
\end{array}\right)$$
可逆.

\;\\
九，是否存在$\mathrm{A},\mathrm{B}\in \mathrm{Z}^{n\times n}$，使得
$$\mathrm{det}(\mathrm{A}^2+\mathrm{B}^2)=2023$$

\;\\
十，设$\mathrm{A},\mathrm{B}$是方阵，$\mathrm{A}^{100}=0,\mathrm{B}=\mathrm{A}\mathrm{B}-\mathrm{B}\mathrm{A}$,证明：$\mathrm{B}=0$

\;\\
十一,(附加题)试找出所有的满射$f:\mathrm{F}^{n\times n}\rightarrow \{0,1\dots n\}$,满足$f(\mathrm{A}\mathrm{B})\leq\mathrm{min}\{f(\mathrm{A}),f(\mathrm{B})\}$.





\end{document}